% Abhakseite „Das kann ich“ – Zone nur im Gesamt (\mitzone)
%% TODO Ich-kann-Sätze: Platzhalter wie in den Aufgabendateien
\begin{abhakseite}
\ifdefined\mitzone
\abhakgruppe{Kennst du schon}
\abhak{1}{Verteilungstabelle einer Zufallsgröße aufstellen und lesen}
\abhak{2}{Pfadregeln und Gegenereignis (mit und ohne Zurücklegen)}
\abhak{3}{Bernoulli-Kette erkennen}
\abhak{4}{Brüche gewichten und addieren, Prozentrechnung}
\abhak{5}{Lineare und quadratische Gleichungen sowie einfache Ungleichungen lösen, Lösungen sieben}
\abhak{6}{Wurzelterme bilden und vergleichen}
\abhak{7}{Pfadregeln und Gegenereignis (mit und ohne Zurücklegen) – Fehler finden}
\abhak{8}{Pfadregeln und Gegenereignis (mit und ohne Zurücklegen) – selbst rechnen}
\fi
\abhakgruppe{Einheit 1 · Erwartungswert berechnen und deuten: die Verteilung beschaffen (Tabelle, Sachtext, Baumpfade, Kosten je Ausgang; fehlende Wahrscheinlichkeit über die Summe 1), die gewichtete Summe bilden (bei Anzahlen n · p), das Ergebnis deuten}
\abhak{9}{Erwartungswert berechnen und deuten}
\abhak{10}{Prozentuale Abweichung einer Anzahl vom Erwartungswert berechnen}
\abhak{11}{Erwartungswert der Augensumme aus dem Erwartungswert der Trefferzahl berechnen}
\abhak{12}{Gesamtzahl der Gruppen aus der Personenzahl und der mittleren Gruppengröße unter einer vereinfachenden Annahme schätzen}
\abhak{13}{Erwartungswert berechnen und deuten – Fehler finden, Begründen, Darstellungswechsel, Anwendung}
\abhakgruppe{Einheit 2 · Rückwärts}
\abhak{14}{Rückwärts}
\abhak{15}{Erwartungswertgleichung für einen Glücksradparameter aus den Spielregeln herleiten}
\abhak{16}{Zwei Wahrscheinlichkeiten einer Verteilung aus dem Erwartungswert und der Summe 1 bestimmen}
\abhak{17}{Restwahrscheinlichkeit und Erwartungswert eines Teilgewinns aus dem Gesamterwartungswert bestimmen}
\abhak{18}{Verhältnis zweier Auszahlungen aus dem Ausgleich der Erwartungswerte berechnen}
\abhak{19}{Kugelzahl für den größten Erwartungswert der Auszahlung ermitteln}
\abhak{20}{Untere Schranke für eine Kugelbeschriftung aus dem Erwartungswert der Summe ohne Rechnung begründen}
\abhak{21}{Aussage über gleiche Erwartungswerte aufeinanderfolgender Parameterwerte über eine Gleichung beurteilen}
\abhak{22}{Rückwärts – Fehler finden, Begründen, Anwendung}
\abhakgruppe{Einheit 3 · Einheit 3}
\abhak{23}{Streuung}
\abhak{24}{Stichprobenumfang aus einer vorgegebenen Standardabweichung der Binomialverteilung berechnen}
\abhak{25}{Verhältnis der Varianzen zweier Binomialverteilungen aus dem Erwartungswert im Diagramm bestimmen}
\abhak{26}{Streuung – Fehler finden, Begründen, Darstellungswechsel, Anwendung}
\abhakgruppe{Einheit 4 · Einheit 4}
\abhak{27}{Diagramm}
\abhak{28}{Diagramm – Fehler finden, Begründen, Darstellungswechsel, Anwendung}
\end{abhakseite}
